% ===========================================================================
\section{Conclusion}
\label{sec:conclusion}
% ===========================================================================

We presented \system{}, a symptom-guided fault reproduction orchestrator that
synthesizes minimal fault recipes for distributed systems. Through greedy
dimensional reduction with binary search, \system{} converges on the boundary
between safe and dangerous operation in $O(\log_2 n)$ iterations per
dimension, using 52\% fewer trials than exhaustive grid approaches while
producing tighter, more actionable boundary specifications.

Our comprehensive evaluation across 10 production systems and 5 fault types
reveals three important properties of the distributed systems ecosystem:

\begin{enumerate}[leftmargin=*,topsep=4pt,itemsep=2pt]
  \item \textbf{Consensus systems are universally fragile.} All Raft-based
        systems exhibit catastrophic failure boundaries at sub-20\,ms network
        delay, a regime entirely below production monitoring visibility.
  \item \textbf{Fault type sensitivity varies by 4 to 16$\times$.} Network faults
        are the most dangerous category; filesystem, CPU, and memory faults
        trigger at proportionally higher severities due to more generous
        timeout configurations.
  \item \textbf{Multi-fault interactions lower boundaries by $\sim$30\%.}
        Single-fault testing produces optimistic boundary estimates; combined
        faults reveal tighter danger zones that better approximate production
        conditions.
\end{enumerate}

These findings have immediate practical implications: monitoring systems should
operate at 10\,ms granularity for consensus-bearing network links (250$\times$
finer than current defaults), CI pipelines should verify resilience boundaries
on every change to consensus code, and system designers should consider adaptive
failure detection to widen danger zones without sacrificing detection latency.

\system{}'s minimal recipes provide a concrete, executable specification of
each system's vulnerability surface, suitable for CI regression tests, SLA
definitions, and architectural decision-making. We release \system{} as open
source to enable the community to characterize and improve the resilience of
distributed systems.

\para{Future work} Three directions merit further investigation:
(1)~\emph{correlated multi-node faults} that model rack-level or AZ-level
degradation simultaneously affecting multiple system participants;
(2)~\emph{production workload replay} to characterize boundaries under
realistic access patterns rather than synthetic benchmarks; and
(3)~\emph{adaptive oracle construction} using machine learning to
automatically generate symptom specifications from production incident reports.
